\documentclass[a4paper,fontsize=12pt,fleqn]{scrartcl}%fleqn pour aligner les equations à gauche

\usepackage{../../Styles/Eval}

\newcommand{\chnum}{Chapitre 10}
\newcommand{\chtitle}{Polarité des molécules et titrages colorimétriques}
\newcommand{\dtype}{DS5}

\begin{document}
\maketitle

\section{La calcite}
C'est après avoir fait tomber au sol un morceau de calcite, et intrigué par la régularité de la fracture, que René Just Haüy (1743-1822) se lance dans l'observation et l'étude des cristaux. Ses études le conduisent à conclure que les cristaux sont des composés organisés et il introduit alors la notion de molécule intégrante remplacée plus tard par le terme maille cristalline.\\
La calcite pure est un solide ionique translucide et blanchâtre constitué d'ions calcium $Ca^{2+}$ et d'ions carbonate $CO_3^{2-}$ lorsque les eaux de pluie s'écoulent et ruissellent sur les sols riches en calcite, la calcite se dissout dans l'eau qui s'enrichit alors en ions calcium.

\begin{enumerate}
    \item Quelles sont les interactions à l'origine de la cohésion de la calcite ? Donner une explication de ces interactions. (3pts)
    \item Écrire l'équation de dissolution de la calcite dans l'eau. (1pt)
    \item Calculer la concentration maximale, exprimée en \unit{\gram\per\liter}, des ions calcium $Ca^{2+}$ et carbonate $CO_3^{2-}$ dans un litre d'eau douce. (3pts)
\end{enumerate}
\begin{data}
    \begin{itemize}
        \item \textbf{Solubilité de la calcite :} $s = \SI{50}{ppm}$ environ à \SI{20}{\celsius}.
        \item \textbf{Masse volumique de la solution :} $\rho = \SI{1,0}{\gram\per\centi\meter\cubed}$
        \item \SI{1}{ppm} signifie \SI{1}{g} dans \SI{1000000}{g} de solution
    \end{itemize}
\end{data}

\newpage 
\section{Extraction du limonène}
\begin{minipage}{.49\linewidth}
    Le limonène est la molécule responsable de l'arôme d'orange, sa formule est donnée ci-contre. Lorsque l'on fabrique une huile essentielle d'orange, une petite partie du liminène se retrouve en solution aqueuse, et pas dans l'huile essentielle. On souhaite extraire ce liminène résiduel.
\end{minipage}
\begin{minipage}{.49\linewidth}
    \centering 
    \chemfig{
        CH_3-[6]C*6(-CH_2-CH_2-CH(-[6]C(=[5]H_2C)-[7]CH_3)-CH_2-CH=)
        }
\end{minipage}

\begin{enumerate}
    \item Solubilité \begin{enumerate}
        \item Le limonène est-il polaire ? (2pts)
        \item Justifier qu'il soit très peu soluble dans l'eau, mais très soluble dans le cyclohexane. (2pts)
    \end{enumerate}
    \item Miscibilité \begin{enumerate}
        \item Quels types de liaisons inter-moléculaires peuvent exister entre l'eau et l'éthanol ? En faire une représentation schématique. (2pts)
        \item Justifier que l'eau et l'éthanol sont miscibles. (2pts)
    \end{enumerate}
    \item Quel solvant doit-on utiliser pour réaliser cette extraction ? Faire un schéma légendé de l'ampoule à décanter après agitation et repos. (3pts)
    \item Pour séparer le limonène du solvant extracteur, on veut évaporer le solvant sans évaporer le limonène. À quelle température soit-on se placer ? (2pts)
\end{enumerate}
\begin{table}[h!]
    \centering 
    \begin{tabular}{|c|c|c|c|c|}
        \hline
        \textbf{Nom} & \textbf{Formule} & $\theta_{eb}$ (\unit{\celsius}) & \textbf{Solubilité du limonène} & \textbf{Miscibilité avec l'eau}\\
        \hline
        Limonène & $C_{10}H_{16}$ & 176 & & \\
        \hline
        Eau & $H_2O$ 100 & Très faible & & \\
        \hline
        Cyclohexane & $C_6H_{12}$ & 81 & Bonne & Non \\
        \hline
        Éthanol & $CH_3CH_2OH$ & 79 & Bonne & Oui \\
        \hline
    \end{tabular}
\end{table}

\end{document}

